\section{Additional Proofs}
\label{sec:proofs}

\subsection{Proof of Lemma~\ref{lem:ML2-bound}}
\label{subsec:ML2-bound}

The exact integer objects in Algorithm~\ref{alg:ML2} are the vectors $b_i$, or
equivalently the exact integral Gram matrix $G$. The quantities
$\bar r_{i,j}$, $\bar\mu_{i,j}$, $\bar s^{(\kappa)}_j$, and $\bar\delta$ are
fixed-precision floating-point numbers, and the indices $\kappa,\zeta$ do not
affect the integer-size bound.

We check the only steps that can affect exact integer data.

\smallskip
\noindent\textbf{(1) Initialization.}
In line~1, the algorithm computes
\[
  G_{i,j}=\langle b_i,b_j\rangle .
\]
By Cauchy--Schwarz,
\[
  |G_{i,j}|
  \le \|b_i\|\,\|b_j\|.
\]

\smallskip
\noindent\textbf{(2) Lazy size-reduction.}
Line~4 calls Algorithm~\ref{alg:lazy-size-reduction}. By
Lemma~\ref{lem:lazy-size-reduction-bound}, this subroutine does not increase
the size of any exact integer object beyond $B_{\rm in}$.

\smallskip
\noindent\textbf{(3) Insertion and remaining updates.}
Lines~5--12 only compare or copy floating-point quantities. Line~13 inserts
$b_{\kappa'}$ right before $b_\kappa$ and updates $G$ accordingly; this is just
a permutation of the current vectors and of the corresponding rows and columns
of $G$, so it creates no new integer magnitude. Lines~14--17 only test whether
$b_\kappa=0$ and update the indices $\zeta$ and $\kappa$.

Therefore every iteration of Algorithm~\ref{alg:ML2} preserves the invariant
that all exact integer objects are bounded by $B_{\rm in}$. Since the invariant
holds after initialization, it holds throughout the algorithm. Hence the exact
integer size in Algorithm~\ref{alg:ML2} is bounded by
\[
  \max_{1\le i\le d}\|b_i\|^2.
\] \qed

\begin{algorithm}
\caption{\textsf{LazySizeReduce}$(\kappa,\zeta,(b_1,\dots,b_d),G,\bar r,\bar\mu)$~\cite[Figure~5]{NS09}}
\label{alg:lazy-size-reduction}
\begin{algorithmic}[1]
  \Statex \textbf{Input:} An index $\kappa$, vectors
  $(b_{\zeta+1},\dots,b_\kappa)$, their Gram matrix $G$, floating-point numbers
  $\bar r_{i,j}$ and $\bar\mu_{i,j}$ for $\zeta+1\le j\le i<\kappa$,
  and fp-precision $\ell+1=53$.
  \Statex \textbf{Output:} Updated $b_\kappa$, $G$, and floating-point numbers
  $\bar r_{\kappa,j}$, $\bar\mu_{\kappa,j}$, and
  $\bar s^{(\kappa)}_j$ for $\zeta+1\le j\le \kappa$.

  \State $\bar\eta \gets \diamond(\diamond(0.55)+1/2)/2$.
  \State Compute the $\bar r_{\kappa,j}$'s, $\bar\mu_{\kappa,j}$'s, and
  $\bar s^{(\kappa)}_j$'s with steps 2--7 of the CFA with ``$i=\kappa$''
  for the vectors $b_{\zeta+1},\dots,b_\kappa$.
  \If{$\max_{\zeta+1\le j<\kappa} |\bar\mu_{\kappa,j}| \le \bar\eta$}
    \State \Return $(b_\kappa,G,\bar r_{\kappa,*},\bar\mu_{\kappa,*},\bar s^{(\kappa)}_*)$.
  \Else
    \For{$i=\kappa-1$ \textbf{down to} $\zeta+1$}
      \State $X_i \gets \lfloor \bar\mu_{\kappa,i}\rceil$.
      \For{$j=\zeta+1$ \textbf{to} $i-1$}
        \State $\bar\mu_{\kappa,j}\gets
        \diamond\!\left(\bar\mu_{\kappa,j}
        -\diamond(X_i\bar\mu_{i,j})\right)$.
      \EndFor
    \EndFor
    \State $b_\kappa \gets b_\kappa-\sum_{i=\zeta+1}^{\kappa-1}X_i b_i$,
    and update $G$ accordingly.
    \State \textbf{go to} line 2.
  \EndIf
\end{algorithmic}
\end{algorithm}

\begin{lemma}[Integer-size bound for \textsf{LazySizeReduce}]
\label{lem:lazy-size-reduction-bound}
Let
\[
  B_{\rm in}:=\max_{1\le i\le d}\|b_i\|^2
\]
be the initial input bound of Algorithm~\ref{alg:ML2}. Assume that before a
call to Algorithm~\ref{alg:lazy-size-reduction}, every exact integer object
currently stored by the algorithm, namely the vector representation or
equivalently the exact Gram-matrix representation, is bounded by $B_{\rm in}$.
Then the returned objects
\[
  b_\kappa,\qquad G,\qquad
  \bar r_{\kappa,*},\quad \bar\mu_{\kappa,*},\quad
  \bar s^{(\kappa)}_*
\]
do not increase the exact integer size beyond $B_{\rm in}$.
\end{lemma}

\begin{proof}
Algorithm~\ref{alg:lazy-size-reduction} manipulates two types of data: exact
integer data, namely $b_\kappa$ and the Gram matrix $G$, and fixed-precision
floating-point data, namely $\bar r$, $\bar\mu$, $\bar s$, and $\bar\eta$.
The floating-point quantities do not contribute to the integer-size bound.

Line~1 only assigns the floating-point value $\bar\eta$. Line~2 computes the
floating-point GSO quantities by the CFA procedure, using $G$ only as exact
input data and producing only rounded floating-point values. Hence lines~1--4
do not create any new exact integer magnitude.

The only nontrivial exact integer update occurs in lines~6--11, where
\[
  X_i\gets \lfloor \bar\mu_{\kappa,i}\rceil,
  \qquad
  b_\kappa\gets b_\kappa-\sum_{i=\zeta+1}^{\kappa-1}X_i b_i,
\]
and the exact Gram matrix $G$ is updated accordingly. This is precisely the
SizeReduce operation analyzed in~\cite{KLKL25}. By the SizeReduce integer-size
bound, this update does not increase the fixed exact-integer budget: after the
update, the reduced vector $b_\kappa$, or equivalently the updated row and
column of $G$, is again represented within the bound $B_{\rm in}$. All other
rows and columns of $G$ are unchanged.

Finally, line~12 only repeats the same procedure. Therefore, by induction over
the repetitions of Algorithm~\ref{alg:lazy-size-reduction}, the subroutine
returns without increasing the exact integer size beyond $B_{\rm in}$.
\end{proof}
% \subsection{Proof of Lemma~\ref{lem:mlll-bound}}
% \label{subsec:mlll-bound}

% We follow the same fixed-precision integer model as in the bound analyses of
% LLL/L2/SizeReduce in~\cite{KLKL25}: all \emph{integer} objects that
% Algorithm~\ref{alg:MLLL} manipulates are stored with a fixed budget.
% (namely $h_i$ and, when the input lattice is integral, $b_i$ or an equivalent integer Gram-matrix representation)
% It is trivial that these objects does not increase the maximum size.
% So the only way the \emph{maximum} integer size could increase is if the algorithm
% created an integer outside the fixed representable range before reduction.

% We check the only lines that can create new integers.

% \smallskip
% \noindent\textbf{(1) Lines 2--5.}
% In line 2, computing $\mu_{\beta j}$ and computing $B_\beta$ need to compute inner product of two vectors so the size during this computation is at most \[\max_{1\le i\le g} ||a_i||^2.\]
% Lines 3--5 only assign zeros, copy existing vectors, or shift
% indices. These operations do not create any new integer magnitude, hence do not
% increase the maximum size.

% \smallskip
% \noindent\textbf{(2) Lines 7--12.}
% When $|\mu_{ml}|>1/2$, the algorithm sets $t\gets\lfloor \mu_{ml}\rceil\in\mathbb{Z}$
% and performs
% \[
% b_m \leftarrow b_m - t b_l,\qquad h_m \leftarrow h_m - t h_l,
% \]
% together with the standard updates of $\mu_{ml},\mu_{mj}$.
% Here $t = \lfloor\mu_{ml}\rceil$, so the maximum size does not grow by the SizeReduce analysis in~\cite{KLKL25}.
% After the updates, all integer entries of $b_m$ and $h_m$ are reduced back to the fixed representative range, so the maximum integer size cannot increase in this loop.
% (The remaining updates in line~9 are performed on the real-valued $\mu$'s and
% do not introduce any new unbounded integer.)

% \smallskip
% \noindent\textbf{(3) Lines 16--26.}
% The assignments in lines~19--21 and lines~24--26 are permutations of existing objects so it is trivial that the size does not grow.
% In line 16, since $0 < \mu \le 1$, the size does not grow.\\
% In line 23, since $\sigma < g$ by line~25, it is trivial that the size does not grow.

% \smallskip
% Combining (1)--(3), we conclude the stated bound in Lemma~\ref{lem:mlll-bound} for the size of integers in modified LLL algorithm.\\ \qed


\subsection{Proof of Lemma~\ref{lem:nrd-mod}}
\label{subsec:nrd-mod}
By multiplicativity of the reduced norm,
\[
\nrd(\alpha)
=
\nrd(\gamma\beta)
=
\nrd(\gamma)\nrd(\beta).
\]
Since $\gcd(\nrd(\beta),N)=1$, we also have
$\gcd(\nrd(\beta),N^2)=1$.
Hence multiplication by $\nrd(\beta)$ does not change the gcd with $N^2$, and therefore
\[
\gcd\!\bigl(N^2,\nrd(\alpha)\bigr)
=
\gcd\!\bigl(N^2,\nrd(\gamma)\nrd(\beta)\bigr)
=
\gcd\!\bigl(N^2,\nrd(\gamma)\bigr)
=
N.
\]

Now write
\[
\widetilde{\alpha}=\alpha+N\delta
\qquad\text{for some }\delta\in\mathcal O_0.
\]
Using $\nrd(x)=x\bar x$, we obtain
\begin{align*}
\nrd(\widetilde{\alpha})
&=
(\alpha+N\delta)(\bar\alpha+N\bar\delta) \\
&=
\alpha\bar\alpha
+
N(\alpha\bar\delta+\delta\bar\alpha)
+
N^2\delta\bar\delta \\
&=
\nrd(\alpha)
+
N\,\operatorname{tr}(\alpha\bar\delta)
+
N^2\nrd(\delta).
\end{align*}
Because $\delta\in\mathcal O_0$, we have
\[
\bar\delta=\operatorname{tr}(\delta)-\delta\in\mathcal O_0,
\]
hence $\alpha\bar\delta\in\mathcal O_0$, and therefore
\[
\operatorname{tr}(\alpha\bar\delta)\in\mathbb Z.
\]
It follows that
\[
\nrd(\widetilde{\alpha})
\equiv
\nrd(\alpha)
\pmod N.
\]
By Item~1, $N\mid \nrd(\alpha)$, hence also
\[
N\mid \nrd(\widetilde{\alpha}).
\]

Finally, from $\widetilde{\alpha}=\alpha+N\delta$ and $N\delta\in N\mathcal O_0=\mathcal O_0N$, we get
\[
\mathcal O_0\widetilde{\alpha}+\mathcal O_0N
\subseteq
\mathcal O_0\alpha+\mathcal O_0N.
\]
Conversely, since
\[
\alpha=\widetilde{\alpha}-N\delta,
\]
we also have
\[
\mathcal O_0\alpha+\mathcal O_0N
\subseteq
\mathcal O_0\widetilde{\alpha}+\mathcal O_0N.
\]
Therefore
\[
\mathcal O_0\langle \widetilde{\alpha},N\rangle
=
\mathcal O_0\langle \alpha,N\rangle.
\] \qed

\subsection{Proof of Lemma~\ref{lem:RandomIdealGivenNorm-distribution}}
\label{subsec:RandomIdealGivenNorm-distribution}

More generally, if $a,a'\in \mathcal O_0$ satisfy $a-a'\in N\mathcal O_0$, then
\[
\mathcal O_0 a+N\mathcal O_0=\mathcal O_0 a'+N\mathcal O_0.
\]
Indeed, writing $a=a'+N\delta$ with $\delta\in\mathcal O_0$, we get
\[
\mathcal O_0 a \subseteq \mathcal O_0 a' + N\mathcal O_0,
\]
hence
\[
\mathcal O_0 a + N\mathcal O_0 \subseteq \mathcal O_0 a' + N\mathcal O_0.
\]
The reverse inclusion follows symmetrically from $a'=a-N\delta$.

Now take $a=\gamma\beta$ and $a'=\overline{\gamma\beta}$. By construction,
\[
\gamma\beta-\overline{\gamma\beta}\in N\mathcal O_0.
\]
Therefore
\[
\mathcal O_0(\gamma\beta)+N\mathcal O_0
=
\mathcal O_0(\overline{\gamma\beta})+N\mathcal O_0.
\]
Thus, for every realization of the random choices, the modified and the original
algorithms return exactly the same left $\mathcal O_0$-ideal. Consequently their
output distributions are identical.\\ \qed

\subsection{Proof of Lemma~\ref{lem:latticeinter-bound}}
\label{subsec:latticeinter-bound}

Lines~5--6 replace \(M_1,M_2\) by LLL-reduced bases of the same lattices.
Thus \(M_k/r_k\) still generates \(I_k\), and \(|\det M_k|\) is unchanged for
\(k=1,2\).  By Lemma~\ref{lem:detnrd},
\[
  \frac{|\det M_k|}{r_k^4}
  =
  \det(I_k)
  =
  \det(\calO_0)\nrd(I_k)^2,
\]
and hence
\[
  |\det M_k|
  =
  r_k^4\det(\calO_0)\nrd(I_k)^2.
\]

We first consider Lines~7--8.  For each \(k\in\{1,2\}\), the input to
\(\mathsf{LatticeDual}\) is \(M_k/r_k\).  By the corrected adjugate bound,
\[
  \|\operatorname{adj}(M_k)\|_\infty
  \le
  8p\,\frac{|\det M_k|}{r_k\sqrt{\nrd(I_k)}}.
\]
In Line~2 of Algorithm~\ref{alg:LatticeDual}, the numerator of the dual basis is
computed as
\[
  D_k = r_k\operatorname{adj}(M_k)^{\mathsf T},
\]
up to the final gcd normalization.
Hence applying Lemma~\ref{lem:detO0},
\[
  \|D_k\|_\infty
  \le
  8p\,\frac{|\det M_k|}{\sqrt{\nrd(I_k)}}
  =
  2p^2\,r_k^4\nrd(I_k)^{3/2}.
\]
The denominator \(s_k\) produced by \(\mathsf{LatticeDual}\) divides
\(|\det M_k|\), so by Lemma~\ref{lem:detO0}, the calls in Lines~7--8 are bounded by
\[
  \max_{k\in\{1,2\}}
  \left(
    r_k^4\frac{p}{4}\nrd(I_k)^2,\,
    2p^2\,r_k^4\nrd(I_k)^{3/2}
  \right).
\]

Next consider Lines~9--10.  Since \(s_k\mid |\det M_k|\), we have
\[
  s=\operatorname{lcm}(s_1,s_2)
  \le
  s_1s_2
  \le
  |\det M_1|\,|\det M_2|.
\]
Using Lemma~\ref{lem:detnrd} and Lemma~\ref{lem:detO0}, this gives
\[
  s
  \le
  r_1^4r_2^4\frac{p^2}{16}
  \nrd(I_1)^2\nrd(I_2)^2.
\]

Now write the columns of \(D_k\) as \(d_{k,j}\).  Since \(C\) is formed as
\[
  C=
  \left[
    \frac{s}{s_1}D_1\ \middle|\ \frac{s}{s_2}D_2
  \right],
\]
we have
\[
  \frac{s}{s_1}\le s_2\le |\det M_2|,
  \qquad
  \frac{s}{s_2}\le s_1\le |\det M_1|.
\]
Moreover,
\[
  \max_j\|d_{k,j}\|^2
  \le
  4\|D_k\|_\infty^2
  \le
  256p^2\frac{|\det M_k|^2}{\nrd(I_k)}.
\]
Therefore the squared length of every column of \(C\) is bounded by
\[
  256p^2\max_{k\in\{1,2\}}\frac{|\det M_k|^2}{\nrd(I_k)}.
\]
Substituting
\[
  |\det M_k|=r_k^4\det(\calO_0)\nrd(I_k)^2
\]
and applying Lemma~\ref{lem:detO0} gives the bound
\[
  16p^4\max_{k\in\{1,2\}} r_k^8\nrd(I_k)^3.
\]

Line~11 applies \(\mathsf{MLLL}\) to \(C\).  By Lemma~\ref{lem:ML2-bound}, the
integer growth in this step is bounded by the squared length of the largest input
column, and hence by the coarse bound above.  This is the step where replacing the
HNF computation by MLLL improves the intermediate-size estimate.

Finally, Line~12 calls \(\mathsf{LatticeDual}\) on \(H/s\), which generates
\[
  I_\Sigma=I_1^\# + I_2^\#.
\]
The corrected adjugate bound gives
\[
  \|\operatorname{adj}(H)\|_\infty
  \le
  8p\,\frac{|\det H|}{s\sqrt{\nrd(I_\Sigma)}}.
\]
Since Algorithm~\ref{alg:LatticeDual} multiplies the adjugate by \(s\), the numerator
computed in this call is bounded by
\[
  8p\,\frac{|\det H|}{\sqrt{\nrd(I_\Sigma)}}.
\]
The denominator is bounded by \(|\det H|\).  Hence Line~12 is bounded by
\[
  \max\left(
    |\det H|,
    8p\,\frac{|\det H|}{\sqrt{\nrd(I_\Sigma)}}
  \right).
\]
By Lemma~\ref{lem:normintersection}, Lemma~\ref{lem:detnrd}, and Lemma~\ref{lem:detO0}, we have the bound
\[
    \max\left(
    \frac{p}{4}\nrd(I_1)^2\nrd(I_2)^2,
    2p^2\,\nrd(I_1)\nrd(I_2)
  \right).
\]

The remaining lines only perform gcd normalization, coordinate reconstruction, or
reduction on already bounded data, and do not introduce larger integers in this
accounting.  Taking the maximum of the contributions above gives the claimed bound
\(B_{\cap}\). \\ \qed


\begin{algorithm}
\caption{LatticeDual($M,r$)}
\label{alg:LatticeDual}
\begin{algorithmic}[1]
  \Statex \textbf{Input:} A full-rank matrix $L = M/r$ representing a $\frac34$-LLL-reduced basis of a left $\calO_0$-ideal $I$, where $M\in\mathbb{Z}^{4\times 4}$ and $r>0$.
  % Here the columns of $L$ are coordinates in the basis $(1,i,j,k)$.
  \Statex \textbf{Output:} A numerator-denominator pair $(M^\#,r^\#)$
  such that the columns of $M^\#/r^\#$ generate $L^\#$.

  \State $(A,\Delta) \gets \mathsf{AdjugateWithDet}(M)$
    \Statex \hspace{\algorithmicindent} \Comment{$M^{-1}=A/\Delta$, where $\Delta=\det(M)$}
  \State $M^\# \gets r A^{\mathsf T}$
  \State $r^\# \gets \Delta$

  \If{$r^\# < 0$}
    \State $M^\# \gets -M^\#$
    \State $r^\# \gets -r^\#$
  \EndIf

  \State $g \gets \gcd\!\left(r^\#,\{M^\#_{a,b}:1\leq a,b\leq4\}\right)$
  \State $M^\# \gets M^\#/g,\quad r^\# \gets r^\#/g$
  \State \textbf{return} $(M^\#,r^\#)$
\end{algorithmic}
\end{algorithm}

\begin{lemma}
\label{lem:dual-bound}
Assume that L is an LLL-reduced basis of a nonzero integral left \(\mathcal O_0\)-ideal \(I\).
Then, during the execution of Algorithm~\ref{alg:LatticeDual}, the maximum size
of integers is at most
\[
  \max\left(
    |\det M|,
    8p\,\frac{|\det M|}{r\cdot\sqrt{\nrd(I)}}
  \right).
\]
\end{lemma}

\begin{proof}
Let
\[
  (A,\Delta)=\mathsf{AdjugateWithDet}(M),
  \qquad
  A=\operatorname{adj}(M),
  \qquad
  \Delta=\det(M).
\]
By the assumed bound for Algorithm~\ref{alg:AdjugateWithDet}, all integers
appearing in Line~1 are bounded by
\[
  \max\left(
    |\Delta|,
    \frac{8p}{r}\,
    \frac{|\Delta|}{\sqrt{\nrd(rI)}}
  \right).
\]
Moreover,
\[
  \|A\|_\infty
  \le
  \frac{8p}{r}\,
  \frac{|\Delta|}{\sqrt{\nrd(rI)}} .
\]

In Line~2, the algorithm computes
\[
  M^\#=rA^{\mathsf T}.
\]
Therefore
\[
  \|M^\#\|_\infty
  =
  r\|A\|_\infty
  \le
  8p\,\frac{|\Delta|}{\sqrt{\nrd(rI)}} .
\]
Line~3 stores
\[
  r^\#=\Delta,
\]
whose size is \(|\Delta|=|\det M|\).

The sign correction in Lines~4--7 does not change absolute values.  In Line~8,
the gcd is computed from \(r^\#\) and the entries of \(M^\#\), so no integer larger
than
\[
  \max\left(
    |\Delta|,
    8p\,\frac{|\Delta|}{\sqrt{\nrd(rI)}}
  \right)
\]
is introduced.  Finally, Line~9 divides \(M^\#\) and \(r^\#\) by this gcd, which can
only decrease their absolute values.

Thus every integer appearing during Algorithm~\ref{alg:LatticeDual} is bounded by
\[
  \max\left(
    |\det M|,
    8p\,\frac{|\det M|}{r\cdot \sqrt{\nrd(I)}}
  \right),
\]
as claimed.
\end{proof}

\begin{algorithm}
\caption{AdjugateWithDet($M$)}
\label{alg:AdjugateWithDet}
\begin{algorithmic}[1]
  \Statex \textbf{Input:} A matrix $M\in\Z^{4\times4}$ representing a $\frac34$-LLL-reduced basis of a left $\calO_0$-ideal $I$.
  \Statex \textbf{Output:} A pair $(A,\Delta)$ such that
  $A=\operatorname{adj}(M)$, $\Delta=\det(M)$, and
  $M^{-1}=A/\Delta$.

  \State $A \gets 0_{4\times 4}$

  \For{$i=1$ to $4$}
    \For{$j=1$ to $4$}
      \State $(s_1,s_2,s_3)\gets\mathsf{SORT}(\{1,2,3,4\}\setminus\{i\})$
      \State $(c_1,c_2,c_3)\gets\mathsf{SORT}(\{1,2,3,4\}\setminus\{j\})$
      \State Let $N^{(i,j)}=(n_{a,b})_{1\le a,b\le 3}$ be the $3\times 3$ matrix
      defined by
      \[
        n_{a,b}=M_{s_a,s_b}
        \qquad (1\le a,b\le 3).
      \]
      \State $\mu_{i,j} \gets
      n_{1,1}(n_{2,2}n_{3,3}-n_{2,3}n_{3,2})
      -n_{1,2}(n_{2,1}n_{3,3}-n_{2,3}n_{3,1})
      +n_{1,3}(n_{2,1}n_{3,2}-n_{2,2}n_{3,1})$
      \State $A_{j,i} \gets (-1)^{i+j}\mu_{i,j}$
    \EndFor
  \EndFor

  \State $\Delta \gets
  m_{1,1}A_{1,1}
  +m_{1,2}A_{2,1}
  +m_{1,3}A_{3,1}
  +m_{1,4}A_{4,1}$

  \If{$\Delta=0$}
    \State \textbf{error} ``$M$ is not full rank''
  \EndIf

  \State \textbf{return} $(A,\Delta)$
\end{algorithmic}
\end{algorithm}

\begin{lemma}
\label{lem:adjugate-bound}
During the execution of Algorithm~\ref{alg:AdjugateWithDet}, the maximum size of integers is less than \[\max\left(|\det M|,\,\frac{8p}{r}\,\frac{|\det M|}{\sqrt{\nrd(I)}}\right).\]
\end{lemma}

\begin{proof}
It is clear that the maximum size of integers appears in line~11.
Since $\Delta = \det(M)$, by Lemma~\ref{lem:adjugate-lll-O0-bound}, we have the maximum size of integers
\[\max\left(|\det M|,\,p\,\frac{|\det(M)|}{\sqrt{\nrd(I)}}\right).\]
\end{proof}

\begin{lemma}[Adjugate bound for LLL-reduced \(\mathcal O_0\)-ideal bases]
\label{lem:adjugate-lll-O0-bound}
Let $G_0$ be the Gram matrix of $\nrd$ in the basis $(1,i,j,k)$.
Set
\[
  C_0
  :=
  8\,\frac{\sqrt{\det(G_0)}}{\lambda_0^{3/2}}
  =
  8p.
\]

Let \(I\subseteq \mathcal O_0\) be a nonzero proper integral left
\(\mathcal O_0\)-ideal, and let
$(\alpha_1,\ldots,\alpha_4)$ be a \(3/4\)-LLL-reduced basis of \(I\).
Let
\[
  \frac{M}{r}\in\mathbb Q^{4\times 4},
  \qquad M\in\mathbb Z^{4\times 4},\quad r>0,
\]
be the coordinate matrix; that is,
the \(\ell\)-th column of \(M/r\) is the \((1,i,j,k)\)-coordinate vector of \(\alpha_\ell\).
Then, the stored entries of \(\operatorname{adj}(M)\) and \(\det(M)\) are
bounded by
\[
  \max\left(
    |\det M|,\,
    \frac{8p}{r}\,\frac{|\det M|}{\sqrt{\nrd(I)}}
  \right).
\]
\end{lemma}

\begin{proof}
Let \(\|\cdot\|_{\nrd}\) be the Euclidean norm induced by \(\nrd\), so that
\[
  \|\alpha\|_{\nrd}^2=\nrd(\alpha).
\]
If \(m/r\in\mathbb R^4\) is the \((1,i,j,k)\)-coordinate vector of \(\alpha\), then
\[
  \|\alpha\|_{\nrd}^2
  =
  (m/r)^{\mathsf T}G_0(m/r).
\]
Since \(\lambda_0=\lambda_{\min}(G_0)=1\), we have
\[
  \|m\|_2
  \le
  r\,\|\alpha\|_{\nrd}.
\]

Let \(m_1,\ldots,m_4\in\mathbb Z^4\) be the columns of \(M\), corresponding to
\(\alpha_1,\ldots,\alpha_4\).  The \((j,i)\)-entry of
\(\operatorname{adj}(M)\) is, up to sign, the determinant of the \(3\times 3\)
minor obtained by deleting row \(i\) and column \(j\).  This determinant is the
volume of the projection of the three vectors \(m_\ell\), \(\ell\neq j\), onto a
coordinate hyperplane.  By Hadamard's inequality,
\[
  \left|\operatorname{adj}(M)_{j,i}\right|
  \le
  \prod_{\ell\neq j}\|m_\ell\|_2 .
\]
Using the comparison between the standard Euclidean norm on numerator
coordinates and \(\|\cdot\|_{\nrd}\), we obtain
\[
  \left|\operatorname{adj}(M)_{j,i}\right|
  \le
  r^3
  \prod_{\ell\neq j}\|\alpha_\ell\|_{\nrd}.
\]

Since \(\alpha_1,\ldots,\alpha_4\) is LLL-reduced, its orthogonality
defect is bounded by
\[
  \prod_{\ell=1}^4 \|\alpha_\ell\|_{\nrd}
  \le
  2^{4\cdot 3/4}\operatorname{covol}_{\nrd}(I)
  =
  8\,\operatorname{covol}_{\nrd}(I).
\]
Moreover,
\[
  \operatorname{covol}_{\nrd}(I)
  =
  \sqrt{\det\!\left((M/r)^{\mathsf T}G_0(M/r)\right)}
  =
  \frac{|\det M|}{r^4}\sqrt{\det(G_0)}
  =
  \frac{p|\det M|}{r^4}.
\]

We now use the ideal structure.  Since the norm ideal \(\nrd(I)\) is generated by
the reduced norms of elements of \(I\), every nonzero \(\alpha\in I\) satisfies
\[
  \nrd(\alpha)\in \nrd(I)\mathbb Z_{>0}.
\]
As \(B_{p,\infty}\) is definite, this implies
\[
  \|\alpha\|_{\nrd}
  =
  \sqrt{\nrd(\alpha)}
  \ge
  \sqrt{\nrd(I)}.
\]
In particular,
\[
  \prod_{\ell\neq j}\|\alpha_\ell\|_{\nrd}
  =
  \frac{\prod_{\ell=1}^4\|\alpha_\ell\|_{\nrd}}
       {\|\alpha_j\|_{\nrd}}
  \le
  \frac{8\,\operatorname{covol}_{\nrd}(I)}{\sqrt{\nrd(I)}}
  =
  \frac{8p|\det M|}{r^4\sqrt{\nrd(I)}}.
\]
Therefore
\[
  \left|\operatorname{adj}(M)_{j,i}\right|
  \le
  r^3
  \frac{8p|\det M|}{r^4\sqrt{\nrd(I)}}
  =
  \frac{8p}{r}\,
  \frac{|\det M|}{\sqrt{\nrd(I)}}.
\]
Taking the maximum over all \(i,j\) proves the adjugate bound.

The determinant \(\det(M)\) is also stored during
Algorithm~\ref{alg:AdjugateWithDet}, so the stored entries of
\(\operatorname{adj}(M)\) and \(\det(M)\) are bounded by
\[
  \max\left(
    |\det M|,\,
    \frac{8p}{r}\,\frac{|\det M|}{\sqrt{\nrd(I)}}
  \right).
\]
\end{proof}

% \section{Analysis on Minkowski-reduced Basis Computation}
% \label{sec:minkowski}

% An algorithm for computing the Minkowski-reduced basis, given a lattice and its basis is as follows: (Algorithm~\ref{alg:Minkowski})

% \begin{algorithm}[H]
% \caption{Minkowski($b_1,\ldots,b_d$)~\cite[Section~4.1]{NS09}}
% \label{alg:Minkowski}
% \begin{algorithmic}[1]
%   \Statex \textbf{Input:} A basis $(b_1,\ldots,b_d)$ with $d\le4$.
%   % with its Gram matrix $G=(g_{i,j})_{1\le i,j\le d}$ and $d\le4$.
%   \Statex \textbf{Output:} An ordered basis of $L[b_1,\ldots,b_d]$.

%   \If{$d=1$}
%     \State \textbf{return} $b_1$
%   \EndIf
%   \Repeat
%     \State $(b_1,\ldots,b_d) \gets \mathbf{LLL}(b_1,\dots,b_d;\delta=\frac 34)$
%     \State Compute the Gram matrix $G$ of $(b_1,\dots,b_d)$
%     \State $(b_1,\ldots,b_{d-1}) \gets \mathbf{Minkowski}(b_1,\ldots,b_{d-1})$
%     \State $c \gets \mathbf{ExactPrefixCVP}\left((b_1,\dots,b_{d-1}), b_d\right)$
%     % \State Compute a vector $c \in L[b_1,\ldots,b_{d-1}]$ closest to $b_d$
%     \State $b_d \gets b_d - c$ and update the Gram matrix efficiently
%   \Until{$\|b_d\| \ge \|b_{d-1}\|$}
%   \State \textbf{return} $[b_1,\ldots,b_d]_{\le}$
% \end{algorithmic}
% \end{algorithm}

% \begin{lemma}
% \label{lem:minkowski-bound}
% During the execution of Algorithm~\ref{alg:Minkowski} with an input basis $(b_1,\dots,b_d)$ of a lattice $I$ in $B_{p,\infty}$, the size of integers is less than
% \[\|b_d^\ast\|^2+9(d-1)^2\|b^\star_{d-1}\|^2,\]
% where $(b^\ast_1,\dots,b^\ast_d)$ is an LLL-reduced basis and $(b^\star_1,\dots,b^\star_d)$ is a Minkowski-reduced basis of $I$.
% \end{lemma}

% \begin{proof}
% \noindent\textbf{(1) Lines 5--6.}
% In line~5, computing LLL-reduced basis does not increase the maximum size of integers. (e.g. See~\cite[Lemma~21]{KLKL25})
% In line~6, computing Gram matrix, the maximum size of integers is $\|b_d^\ast\|^2$.

% \smallskip
% \noindent\textbf{(2) Lines 7--9.}
% In line~8, the maximum size of integers during the execution of \textbf{ExactPrefixCVP} is less than
% \[ \|b_d^\ast\|^2 + 9(d-1)^2\|b_{d-1}^\star\|^2, \]
% by Lemma~\ref{lem:cvp-bound}.
% In line~9, it is trivial that the size of integers does not grow when reducing the size of $b_d$, by subtracting the closest vector of $b_d$ to the lattice generated by $(b_1,\dots,b_{d_1})$.

% \end{proof}


% \begin{algorithm}[!t]
% \caption{ExactPrefixCVP($B,t$)~\cite[Section~5]{NS09}}
% \label{alg:ExactPrefixCVP}
% \begin{algorithmic}[1]
% \Require $B = (b_1,\dots,b_m)$ is a Minkowski-reduced basis; 
%          $t \in \mathbb{R}^n$; $\|b_i\| < \|t\|$ for all $i$
% \Ensure A vector $c^\star \in L(B)$ minimizing $\|t - c^\star\|$, where
%         $L(B) = \left\{\sum_{i=1}^m z_i b_i : z_i \in \mathbb{Z}\right\}$.
%     \If{$m = 1$}
%         \State $\alpha \gets \dfrac{\langle b_1,t\rangle}{\langle b_1,b_1\rangle}$
%         \State $x^\star \gets \lfloor\alpha\rceil$
%         \State \Return $x^\star b_1$
%     \EndIf

%     \State $\mathbf{G} \gets \bigl(\langle b_i,b_j\rangle\bigr)_{1 \le i,j \le m}$
%     \State $\mathbf{v} \gets \bigl(\langle b_1,t\rangle,\dots,\langle b_m,t\rangle\bigr)^\top$
%     % \State Compute $\widetilde{\mathbf y}$ such that there exists a unique $\mathbf y$ with $\mathbf G\mathbf y=\mathbf v$ and $\|\widetilde{\mathbf y}-\mathbf y\|_\infty \le 1/2$
%     \State $\mathbf y \gets \mathbf{G}^{-1}\mathbf v$, \quad $y_i \gets \mathbf y _i$
%     \State $\mathcal{S} \gets \displaystyle\prod_{i=1}^m
%            \left\{\lfloor y_i \rfloor - 2,\,
%                   \lfloor y_i \rfloor - 1,\,
%                   \lfloor y_i \rfloor,\,
%                   \lfloor y_i \rfloor + 1,\,
%                   \lfloor y_i \rfloor + 2
%            \right\}$

%     \State $d^\star \gets +\infty$, $c^\star \gets 0$, $\mathbf{x}^\star \gets \bot$

%     \ForAll{$\mathbf{x} = (x_1,\dots,x_m) \in \mathcal{S}$}
%         \State $c \gets \sum_{i=1}^m x_i b_i$
%         \State $d \gets \|t - c\|^2$
%         \If{$d < d^\star$}
%             \State $d^\star \gets d$
%             \State $c^\star \gets c$
%             \State $\mathbf{x}^\star \gets \mathbf{x}$
%         \EndIf
%     \EndFor

%     \State \Return $c^\star$
% \end{algorithmic}
% \end{algorithm}

% \begin{lemma}
% \label{lem:cvp-bound}
% During the execution of Algorithm~\ref{alg:ExactPrefixCVP}, the maximum size of integers is less than $\|t\|^2 + 81\cdot\|b_m\|^2$.
% \end{lemma}

% \begin{proof}
% Let $T := \|t\|$.

% \smallskip
% \noindent\textbf{(1) Lines 1--5.}
% If $m=1$, then
% \[
% \alpha=\frac{\langle b_1,t\rangle}{\langle b_1,b_1\rangle},
% \qquad
% |\langle b_1,t\rangle| \le \|b_1\|\,\|t\| < T^2,
% \qquad
% 0<\langle b_1,b_1\rangle=\|b_1\|^2<T^2.
% \]
% Thus the numerator and denominator occurring in $\alpha$ are bounded by $T^2$. Moreover,
% \[
% |\alpha| \le \frac{\|b_1\|\,\|t\|}{\|b_1\|^2}=\frac{T}{\|b_1\|},
% \]
% so the nearest integer $x^\star$ satisfies
% \[
% |x^\star| \le |\alpha|+\tfrac12 = \frac{T}{\|b_1\|}+\frac12 \le T+\frac12,
% \]
% since we treat only integer vectors. Hence
% \[
% \|x^\star b_1\| \le (|\alpha|+\tfrac12)\|b_1\|
% = T+\frac12\|b_1\|.
% \]

% \smallskip
% \noindent\textbf{(2) Lines 6--7.}
% Assume now that $m\ge 2$.
% In lines 6--7, we form the Gram matrix
% \[
% \mathbf G=(\langle b_i,b_j\rangle)_{1\le i,j\le m}
% \qquad\text{and}\qquad
% \mathbf v=(\langle b_1,t\rangle,\ldots,\langle b_m,t\rangle)^{\top}. 
% \]
% By Cauchy--Schwarz and the hypothesis $\|b_i\|<T$ for every $i$,
% \[
% |G_{ij}| = |\langle b_i,b_j\rangle|
%  \le \|b_i\|\,\|b_j\| < T^2,
% \qquad
% |v_i| = |\langle b_i,t\rangle|
%  \le \|b_i\|\,T < T^2.
% \]
% So all integers appearing in $\mathbf G$ and $\mathbf v$ are less than $T^2$.

% \smallskip
% \noindent\textbf{(3) Line 8.}
% Because $B$ is Minkowski-reduced and the dimension $m$ satisfies $m \le 4$, by Minkowski's second theorem, we have
% \[
% |y_i|\le \|d_i\|\,T
% \le \frac{32}{\pi^2}\cdot\frac{T}{\|b_i\|}.
% \]
% During the compuataion of $\mathbf{G}^{-1}$, since $G$ is the Gram matrix, the maximum size of integers, which occurs by the computation of determinant is less than
% \[ \prod_{i=1}^d \|b_i\|^2 \le \frac{32p^2}{\pi^4}\nrd(I)^2, \]
% by Minkowski's second theorem, where $(b_1,b_2,b_3,b_4)$ is a basis of $I$.

% \smallskip
% \noindent\textbf{(4) Line 12.}
% Fix such an $\mathbf x$, and write $c=B\mathbf x$.
% Because $\mathbf x\in\mathcal S$, for each $i$ we have
% \[
% x_i\in \{\lfloor y_i\rfloor-2,\lfloor y_i\rfloor-1,\lfloor y_i\rfloor,
% \lfloor y_i\rfloor+1,\lfloor y_i\rfloor+2\},
% \]
% Hence, $\|x_i- y_i\| \le 3$, for each $i$.
% By $|x_i|\le |y_i|+2$ we have
% \[
% |x_i|\le \frac{32}{\pi^2}\cdot\frac{T}{\|b_i\|}+2 \le \frac{32}{\pi^2}T+2,
% \]
% since $L(B)$ is integral so $\|b_i\| \ge 1$ and
% It follows that
% \[
% B(\mathbf x-\mathbf y)=\sum_{i=1}^m (x_i-y_i)b_i,
% \]
% so by the triangle inequality,
% \[
% \|B(\mathbf x-\mathbf y)\|
% \le \sum_{i=1}^m |x_i-y_i|\,\|b_i\|
% \le 3\sum_{i=1}^m \|b_i\|
% \le 3m \max_{1\le i\le m}\|b_i\|
% = 3m \|b_3\|.
% \]

% Also, since $\mathbf y=\mathbf G^{-1}\mathbf v$ with $\mathbf G=B^\top B$ and
% $\mathbf v=B^\top t$, we have
% \[
% B^\top(t-B\mathbf y)=\mathbf v-\mathbf G\mathbf y=0.
% \]
% Thus $t-B\mathbf y$ is orthogonal to $\mathrm{span}(B)$, while
% $B(\mathbf y-\mathbf x)\in \mathrm{span}(B)$. Therefore
% \[
% \|t-c\|^2
% = \|t-B\mathbf y\|^2 + \|B(\mathbf y-\mathbf x)\|^2.
% \]
% Since $B\mathbf y$ is the orthogonal projection of $t$ onto $\mathrm{span}(B)$,
% \[
% \|t-B\mathbf y\|\le \|t\|=T,
% \]
% and hence
% \[
% \|t-c\|^2 \le T^2 + (3m\cdot\|b_3\|)^2.
% \]
% Therefore, we have
% \[
% \|c\|\le \|t\|+\|t-c\| \le \bigl(1+\sqrt{1+9m^2}\bigr)T < 14T.
% \]

% \end{proof}


% \begin{algorithm}[H]
% \caption{Inverse3x3($A$)}
% \label{alg:Inverse}
% \begin{algorithmic}[1]
%   \Statex \textbf{Input:} A matrix
%   \[
%     A=
%     \begin{pmatrix}
%       a_{11} & a_{12} & a_{13}\\
%       a_{21} & a_{22} & a_{23}\\
%       a_{31} & a_{32} & a_{33}
%     \end{pmatrix}
%     \in \mathbb{Q}^{3\times 3}.
%   \]
%   \Statex \textbf{Output:} $A^{-1}$ if $A$ is nonsingular.

%   \State $\Delta \gets a_{11}(a_{22}a_{33}-a_{23}a_{32})
%           - a_{12}(a_{21}a_{33}-a_{23}a_{31})
%           + a_{13}(a_{21}a_{32}-a_{22}a_{31})$
%   \If{$\Delta = 0$}
%     \State \textbf{error} ``$A$ is singular''
%   \EndIf

%   \State $C_{11} \gets \phantom{-}(a_{22}a_{33}-a_{23}a_{32})$
%   \State $C_{12} \gets -(a_{21}a_{33}-a_{23}a_{31})$
%   \State $C_{13} \gets \phantom{-}(a_{21}a_{32}-a_{22}a_{31})$

%   \State $C_{21} \gets -(a_{12}a_{33}-a_{13}a_{32})$
%   \State $C_{22} \gets \phantom{-}(a_{11}a_{33}-a_{13}a_{31})$
%   \State $C_{23} \gets -(a_{11}a_{32}-a_{12}a_{31})$

%   \State $C_{31} \gets \phantom{-}(a_{12}a_{23}-a_{13}a_{22})$
%   \State $C_{32} \gets -(a_{11}a_{23}-a_{13}a_{21})$
%   \State $C_{33} \gets \phantom{-}(a_{11}a_{22}-a_{12}a_{21})$

%   \State $A^{-1} \gets \dfrac{1}{\Delta}
%     \begin{pmatrix}
%       C_{11} & C_{21} & C_{31}\\
%       C_{12} & C_{22} & C_{32}\\
%       C_{13} & C_{23} & C_{33}
%     \end{pmatrix}$

%   \State \Return $A^{-1}$
% \end{algorithmic}
% \end{algorithm}